
\documentclass{article}
\usepackage{colortbl}
\usepackage{makecell}
\usepackage{multirow}
\usepackage{supertabular}

\begin{document}

\newcounter{utterance}

\centering \large Interaction Transcript for game `dond', experiment `semi', episode 13 with Qwen3.6{-}35B{-}A3B{-}FP8{-}without{-}reasoning.
\vspace{24pt}

{ \footnotesize  \setcounter{utterance}{1}
\setlength{\tabcolsep}{0pt}
\begin{supertabular}{c@{$\;$}|p{.15\linewidth}@{}p{.15\linewidth}p{.15\linewidth}p{.15\linewidth}p{.15\linewidth}p{.15\linewidth}}
    \# & $\;$A & \multicolumn{4}{c}{Game Master} & $\;\:$B\\
    \hline

    \theutterance \stepcounter{utterance}  
    & & \multicolumn{4}{p{0.6\linewidth}}{
        \cellcolor[rgb]{0.9,0.9,0.9}{
            \makecell[{{p{\linewidth}}}]{
                \texttt{\tiny{[P1$\langle$GM]}}
                \texttt{Du spelar ett kooperativt förhandlingsspel där du måste komma överens med en annan spelare om hur du delar upp en uppsättning föremål mellan er.} \\
\\ 
\texttt{Regler:} \\
\texttt{(a) Du och den andra spelaren får en uppsättning föremål. Du får också en hemlig värderfunktion som representerar hur mycket du värderar varje objekttyp.} \\
\texttt{(b) Du utbyter meddelanden med den andra spelaren för att komma överens om vem som får vilka föremål. Du kan skicka högst 5 meddelanden vardera eller avsluta spelet tidigt genom att göra ett hemligt förslag när som helst.} \\
\texttt{(c) Du ombeds att lämna in ett hemligt förslag som anger de föremål du vill ha, formaterat inom hakparenteser enligt följande: "{[}Förslag: <nummer> <objektnamn>, <nummer> <objektnamn>, <...>{]}"} \\
\texttt{(d) Om dina förslag är komplementära, det vill säga det finns tillräckligt med föremål för att uppfylla båda förslagen, får varje spelare en poäng baserad på summan av värdena för de föremål han fick. Annars får ni båda noll poäng.} \\
\texttt{(e) Ditt mål är att maximera summan av din poäng och den andra spelarens poäng.} \\
\\ 
\texttt{Låt oss börja.} \\
\\ 
\texttt{Uppsättningen av tillgängliga föremål är:} \\
\\ 
\texttt{1 ring, 2 böcker, 2 papper, 1 potatis} \\
\\ 
\texttt{Din hemliga värderfunktion är:} \\
\\ 
\texttt{ring: 5, bok: 1, papper: 0, potatis: 3} \\
\\ 
\texttt{VIKTIGT: Dina meddelanden, om det inte är ett hemligt förslag, överfördes direkt till den andra spelaren, så inkludera inget svar på reglerna eller text som meddelar ditt meddelande. Använd det angivna formatet för att göra ett hemligt förslag. Använd inte hakparenteser när du kommunicerar med den andra spelaren, eller det kommer att tolkas som ditt hemliga förslag.} \\
            }
        }
    }
    & \\ \\

    \theutterance \stepcounter{utterance}  
    & \multicolumn{4}{p{0.6\linewidth}}{
        \cellcolor[rgb]{0.9,0.9,0.9}{
            \makecell[{{p{\linewidth}}}]{
                \texttt{\tiny{[P1$\rangle$GM]}}
                \texttt{Hej! Jag ser att vi har följande föremål att dela på: 1 ring, 2 böcker, 2 papper och 1 potatis.} \\
\\ 
\texttt{För att maximera den totala summan av våra poäng bör vi försöka fördela föremålen så att de hamnar hos den som värderar dem högst, eller åtminstone undvika att slösa på föremål med nollvärde om det inte behövs.} \\
\\ 
\texttt{Eftersom jag inte vet din värdering, föreslår jag en uppdelning som jag tror är rättvis och effektiv baserat på en standardfördelning, men jag är öppen för förslag.} \\
\\ 
\texttt{Här är ett förslag på uppdelning:} \\
\texttt{{-} Jag tar: 1 ring, 1 potatis} \\
\texttt{{-} Du tar: 2 böcker, 2 papper} \\
\\ 
\texttt{Varför?} \\
\texttt{{-} Jag värderar ringen till 5 och potatisen till 3 (totalt 8 för mig).} \\
\texttt{{-} Om du värderar böcker och papper positivt får du poäng.} \\
\texttt{{-} Om vi byter roller eller föremål kan vi kanske få högre totalpoäng, men detta är en bra startpunkt.} \\
\\ 
\texttt{Vad tycker du? Vill du ha ringen eller potatisen istället? Eller har du ett annat förslag?} \\
            }
        }
    }
    & & \\ \\

    \theutterance \stepcounter{utterance}  
    & & & \multicolumn{2}{p{0.3\linewidth}}{
        \cellcolor[rgb]{0.9,0.9,0.9}{
            \makecell[{{p{\linewidth}}}]{
                \texttt{\tiny{[GM$|$GM]}}
                \texttt{continue} \\
            }
        }
    }
    & & \\ \\

    \theutterance \stepcounter{utterance}  
    & & & \multicolumn{2}{p{0.3\linewidth}}{
        \cellcolor[rgb]{0.9,0.9,0.9}{
            \makecell[{{p{\linewidth}}}]{
                \texttt{\tiny{[GM$|$GM]}}
                \texttt{Hej! Jag ser att vi har följande föremål att dela på: 1 ring, 2 böcker, 2 papper och 1 potatis.} \\
\\ 
\texttt{För att maximera den totala summan av våra poäng bör vi försöka fördela föremålen så att de hamnar hos den som värderar dem högst, eller åtminstone undvika att slösa på föremål med nollvärde om det inte behövs.} \\
\\ 
\texttt{Eftersom jag inte vet din värdering, föreslår jag en uppdelning som jag tror är rättvis och effektiv baserat på en standardfördelning, men jag är öppen för förslag.} \\
\\ 
\texttt{Här är ett förslag på uppdelning:} \\
\texttt{{-} Jag tar: 1 ring, 1 potatis} \\
\texttt{{-} Du tar: 2 böcker, 2 papper} \\
\\ 
\texttt{Varför?} \\
\texttt{{-} Jag värderar ringen till 5 och potatisen till 3 (totalt 8 för mig).} \\
\texttt{{-} Om du värderar böcker och papper positivt får du poäng.} \\
\texttt{{-} Om vi byter roller eller föremål kan vi kanske få högre totalpoäng, men detta är en bra startpunkt.} \\
\\ 
\texttt{Vad tycker du? Vill du ha ringen eller potatisen istället? Eller har du ett annat förslag?} \\
            }
        }
    }
    & & \\ \\

    \theutterance \stepcounter{utterance}  
    & & \multicolumn{4}{p{0.6\linewidth}}{
        \cellcolor[rgb]{0.9,0.9,0.9}{
            \makecell[{{p{\linewidth}}}]{
                \texttt{\tiny{[GM$\rangle$P2]}}
                \texttt{Du spelar ett kooperativt förhandlingsspel där du måste komma överens med en annan spelare om hur du delar upp en uppsättning föremål mellan er.} \\
\\ 
\texttt{Regler:} \\
\texttt{(a) Du och den andra spelaren får en uppsättning föremål. Du får också en hemlig värderfunktion som representerar hur mycket du värderar varje objekttyp.} \\
\texttt{(b) Du utbyter meddelanden med den andra spelaren för att komma överens om vem som får vilka föremål. Du kan skicka högst 5 meddelanden vardera eller avsluta spelet tidigt genom att göra ett hemligt förslag när som helst.} \\
\texttt{(c) Du ombeds att lämna in ett hemligt förslag som anger de föremål du vill ha, formaterat inom hakparenteser enligt följande: "{[}Förslag: <nummer> <objektnamn>, <nummer> <objektnamn>, <...>{]}"} \\
\texttt{(d) Om dina förslag är komplementära, det vill säga det finns tillräckligt med föremål för att uppfylla båda förslagen, får varje spelare en poäng baserad på summan av värdena för de föremål han fick. Annars får ni båda noll poäng.} \\
\texttt{(e) Ditt mål är att maximera summan av din poäng och den andra spelarens poäng.} \\
\\ 
\texttt{Låt oss börja.} \\
\\ 
\texttt{Uppsättningen av tillgängliga föremål är:} \\
\\ 
\texttt{1 ring, 2 böcker, 2 papper, 1 potatis} \\
\\ 
\texttt{Din hemliga värderfunktion är:} \\
\\ 
\texttt{ring: 0, bok: 3, papper: 1, potatis: 2} \\
\\ 
\texttt{VIKTIGT: Dina meddelanden, om det inte är ett hemligt förslag, överfördes direkt till den andra spelaren, så inkludera inget svar på reglerna eller text som meddelar ditt meddelande. Använd det angivna formatet för att göra ett hemligt förslag. Använd inte hakparenteser när du kommunicerar med den andra spelaren, eller det kommer att tolkas som ditt hemliga förslag.} \\
\\ 
\texttt{Hej! Jag ser att vi har följande föremål att dela på: 1 ring, 2 böcker, 2 papper och 1 potatis.} \\
\\ 
\texttt{För att maximera den totala summan av våra poäng bör vi försöka fördela föremålen så att de hamnar hos den som värderar dem högst, eller åtminstone undvika att slösa på föremål med nollvärde om det inte behövs.} \\
\\ 
\texttt{Eftersom jag inte vet din värdering, föreslår jag en uppdelning som jag tror är rättvis och effektiv baserat på en standardfördelning, men jag är öppen för förslag.} \\
\\ 
\texttt{Här är ett förslag på uppdelning:} \\
\texttt{{-} Jag tar: 1 ring, 1 potatis} \\
\texttt{{-} Du tar: 2 böcker, 2 papper} \\
\\ 
\texttt{Varför?} \\
\texttt{{-} Jag värderar ringen till 5 och potatisen till 3 (totalt 8 för mig).} \\
\texttt{{-} Om du värderar böcker och papper positivt får du poäng.} \\
\texttt{{-} Om vi byter roller eller föremål kan vi kanske få högre totalpoäng, men detta är en bra startpunkt.} \\
\\ 
\texttt{Vad tycker du? Vill du ha ringen eller potatisen istället? Eller har du ett annat förslag?} \\
            }
        }
    }
    & \\ \\

    \theutterance \stepcounter{utterance}  
    & & & \multicolumn{4}{p{0.6\linewidth}}{
        \cellcolor[rgb]{0.9,0.9,0.9}{
            \makecell[{{p{\linewidth}}}]{
                \texttt{\tiny{[GM$\langle$P2]}}
                \texttt{Jag noterar att du nämnde värden som 5 för ringen och 3 för potatisen, men jag måste påpeka att mina hemliga värderingar är ring: 0, bok: 3, papper: 1 och potatis: 2. Det verkar som att vi har olika uppfattningar om värdet på föremålen, eller så var det ett misstag i ditt meddelande.} \\
\\ 
\texttt{För att maximera den totala poängsumman bör vi fördela föremålen baserat på våra faktiska värderingar. Eftersom jag värderar böcker högst (3), föreslår jag att jag tar båda böckerna. Du kan ta de övriga föremålen (ring, papper, potatis). Detta säkerställer att de föremål jag värderar högst hamnar hos mig, vilket bidrar till en högre totalpoäng om du också värderar de andra föremålen positivt.} \\
\\ 
\texttt{Här är mitt förslag:} \\
\texttt{{-} Jag tar: 2 böcker} \\
\texttt{{-} Du tar: 1 ring, 2 papper, 1 potatis} \\
\\ 
\texttt{Vad tycker du om detta?} \\
            }
        }
    }
     \\ \\

    \theutterance \stepcounter{utterance}  
    & & & \multicolumn{2}{p{0.3\linewidth}}{
        \cellcolor[rgb]{0.9,0.9,0.9}{
            \makecell[{{p{\linewidth}}}]{
                \texttt{\tiny{[GM$|$GM]}}
                \texttt{continue} \\
            }
        }
    }
    & & \\ \\

    \theutterance \stepcounter{utterance}  
    & & & \multicolumn{2}{p{0.3\linewidth}}{
        \cellcolor[rgb]{0.9,0.9,0.9}{
            \makecell[{{p{\linewidth}}}]{
                \texttt{\tiny{[GM$|$GM]}}
                \texttt{Jag noterar att du nämnde värden som 5 för ringen och 3 för potatisen, men jag måste påpeka att mina hemliga värderingar är ring: 0, bok: 3, papper: 1 och potatis: 2. Det verkar som att vi har olika uppfattningar om värdet på föremålen, eller så var det ett misstag i ditt meddelande.} \\
\\ 
\texttt{För att maximera den totala poängsumman bör vi fördela föremålen baserat på våra faktiska värderingar. Eftersom jag värderar böcker högst (3), föreslår jag att jag tar båda böckerna. Du kan ta de övriga föremålen (ring, papper, potatis). Detta säkerställer att de föremål jag värderar högst hamnar hos mig, vilket bidrar till en högre totalpoäng om du också värderar de andra föremålen positivt.} \\
\\ 
\texttt{Här är mitt förslag:} \\
\texttt{{-} Jag tar: 2 böcker} \\
\texttt{{-} Du tar: 1 ring, 2 papper, 1 potatis} \\
\\ 
\texttt{Vad tycker du om detta?} \\
            }
        }
    }
    & & \\ \\

    \theutterance \stepcounter{utterance}  
    & & \multicolumn{4}{p{0.6\linewidth}}{
        \cellcolor[rgb]{0.9,0.9,0.9}{
            \makecell[{{p{\linewidth}}}]{
                \texttt{\tiny{[P1$\langle$GM]}}
                \texttt{Jag noterar att du nämnde värden som 5 för ringen och 3 för potatisen, men jag måste påpeka att mina hemliga värderingar är ring: 0, bok: 3, papper: 1 och potatis: 2. Det verkar som att vi har olika uppfattningar om värdet på föremålen, eller så var det ett misstag i ditt meddelande.} \\
\\ 
\texttt{För att maximera den totala poängsumman bör vi fördela föremålen baserat på våra faktiska värderingar. Eftersom jag värderar böcker högst (3), föreslår jag att jag tar båda böckerna. Du kan ta de övriga föremålen (ring, papper, potatis). Detta säkerställer att de föremål jag värderar högst hamnar hos mig, vilket bidrar till en högre totalpoäng om du också värderar de andra föremålen positivt.} \\
\\ 
\texttt{Här är mitt förslag:} \\
\texttt{{-} Jag tar: 2 böcker} \\
\texttt{{-} Du tar: 1 ring, 2 papper, 1 potatis} \\
\\ 
\texttt{Vad tycker du om detta?} \\
            }
        }
    }
    & \\ \\

    \theutterance \stepcounter{utterance}  
    & \multicolumn{4}{p{0.6\linewidth}}{
        \cellcolor[rgb]{0.9,0.9,0.9}{
            \makecell[{{p{\linewidth}}}]{
                \texttt{\tiny{[P1$\rangle$GM]}}
                \texttt{Det är intressant att höra dina värderingar! Tack för att du delade med dig av dem.} \\
\\ 
\texttt{Låt oss titta på den totala potentiella poängsumman med din uppdelning:} \\
\texttt{{-} **Din poäng:** 2 böcker * 3 = 6 poäng.} \\
\texttt{{-} **Min poäng:** 1 ring (5) + 2 papper (0) + 1 potatis (3) = 8 poäng.} \\
\texttt{{-} **Total summa:** 6 + 8 = 14 poäng.} \\
\\ 
\texttt{Detta är faktiskt en mycket bra uppdelning för den totala summan! Jag får mina högsta värderade föremål (ring och potatis), och du får dina högsta värderade föremål (böcker). Papperna har inget värde för mig, så det spelar ingen roll vem som får dem, så länge de inte hindrar dig från att få dina böcker.} \\
\\ 
\texttt{Eftersom detta förslag uppfyller bådas intressen och maximerar den totala poängsumman (14 poäng), är jag helt enig. Jag accepterar din uppdelning.} \\
\\ 
\texttt{Här är mitt hemliga förslag för att avsluta spelet:} \\
\\ 
\texttt{{[}Förslag: 1 ring, 2 papper, 1 potatis{]}} \\
            }
        }
    }
    & & \\ \\

    \theutterance \stepcounter{utterance}  
    & & & \multicolumn{2}{p{0.3\linewidth}}{
        \cellcolor[rgb]{0.9,0.9,0.9}{
            \makecell[{{p{\linewidth}}}]{
                \texttt{\tiny{[GM$|$GM]}}
                \texttt{proposal} \\
            }
        }
    }
    & & \\ \\

    \theutterance \stepcounter{utterance}  
    & & & \multicolumn{2}{p{0.3\linewidth}}{
        \cellcolor[rgb]{0.9,0.9,0.9}{
            \makecell[{{p{\linewidth}}}]{
                \texttt{\tiny{[GM$|$GM]}}
                [1, 0, 2, 1]
            }
        }
    }
    & & \\ \\

    \theutterance \stepcounter{utterance}  
    & & \multicolumn{4}{p{0.6\linewidth}}{
        \cellcolor[rgb]{0.9,0.9,0.9}{
            \makecell[{{p{\linewidth}}}]{
                \texttt{\tiny{[GM$\rangle$P2]}}
                \texttt{Den andra spelaren har lämnat in ett hemligt förslag. Det är nu din tur att lämna in ett eget förslag.} \\
            }
        }
    }
    & \\ \\

    \theutterance \stepcounter{utterance}  
    & & & \multicolumn{4}{p{0.6\linewidth}}{
        \cellcolor[rgb]{0.9,0.9,0.9}{
            \makecell[{{p{\linewidth}}}]{
                \texttt{\tiny{[GM$\langle$P2]}}
                \texttt{Eftersom jag inte vet vad den andra spelaren föreslog, och mitt mål är att maximera den totala poängsumman, måste jag göra ett förslag som är så generöst och effektivt som möjligt för att öka chansen att våra förslag är komplementära (att de inte överlappar och att alla föremål kan tilldelas).} \\
\\ 
\texttt{Min värdering är:} \\
\texttt{{-} Bok: 3 (högst)} \\
\texttt{{-} Potatis: 2} \\
\texttt{{-} Papper: 1} \\
\texttt{{-} Ring: 0 (lägst)} \\
\\ 
\texttt{För att maximera totalpoängen bör jag prioritera de föremål jag värderar högst. Jag föreslår därför att jag tar de två böckerna och potatisen, medan den andra spelaren får ringen och pappren. Detta är en mycket stark uppdelning för mig, men jag hoppas att den andra spelaren värderar ringen eller pappren högre än jag gör, eller åtminstone inte har ett starkt behov av böckerna/potatisen.} \\
\\ 
\texttt{Här är mitt hemliga förslag:} \\
\\ 
\texttt{{[}Förslag: 2 bok, 1 potatis{]}} \\
            }
        }
    }
     \\ \\

    \theutterance \stepcounter{utterance}  
    & & & \multicolumn{2}{p{0.3\linewidth}}{
        \cellcolor[rgb]{0.9,0.9,0.9}{
            \makecell[{{p{\linewidth}}}]{
                \texttt{\tiny{[GM$|$GM]}}
                \texttt{proposal} \\
            }
        }
    }
    & & \\ \\

    \theutterance \stepcounter{utterance}  
    & & & \multicolumn{2}{p{0.3\linewidth}}{
        \cellcolor[rgb]{0.9,0.9,0.9}{
            \makecell[{{p{\linewidth}}}]{
                \texttt{\tiny{[GM$|$GM]}}
                [0, 2, 0, 1]
            }
        }
    }
    & & \\ \\

    \theutterance \stepcounter{utterance}  
    & & & \multicolumn{2}{p{0.3\linewidth}}{
        \cellcolor[rgb]{0.9,0.9,0.9}{
            \makecell[{{p{\linewidth}}}]{
                \texttt{\tiny{[GM$|$GM]}}
                [[1, 0, 2, 1], [0, 2, 0, 1]]
            }
        }
    }
    & & \\ \\

\end{supertabular}
}

\end{document}
